\chapter{Fixed-rate bond analytics}\label{ch:bonds}
\section{Notation, domain and assumptions}\label{sec:bond-assumptions}
Let \(F>0\) denote face value in currency units, \(c\geq0\) the annual
decimal coupon rate, and \(m\in\{1,2\}\) the number of coupon payments per year.
The nominal annual decimal yield to maturity is \(\yytm\), compounded \(m\)
times per year. Define the dimensionless periodic accumulation factor
\(B=1+\yytm/m>0\).
\begin{table}[htbp]
\centering\small
\caption{Model assumptions and their methodological consequences}\label{tab:assumptions}
\begin{tabular}{@{}L{0.29\textwidth}L{0.64\textwidth}@{}}
\toprule
Assumption & Consequence\\
\midrule
Fixed nonnegative coupons; positive redemption & Cash flows remain fixed in sensitivity calculations.\\
Regular annual/semiannual payments & No irregular first or final coupon periods.\\
Maturity-anchored calendar dates & Invalid month days clip; no business-day or additional month-end rule.\\
Payments strictly after settlement & A settlement-date coupon is excluded; accrual restarts at zero.\\
Actual days within the coupon period & Accrual is a period fraction, not a complete market day-count standard.\\
YTM: fractional coupon periods & Discount at one nominal coupon-compounded annual yield.\\
Zero curve: actual days divided by 365 & Discount separately at continuous annual zero rates.\\
Deterministic instantaneous shocks & Settlement and cash flows remain unchanged.\\
\bottomrule
\end{tabular}
\par\smallskip{\footnotesize Source: conventions in the existing implementation.}
\end{table}
\section{Regular coupon schedule and accrual}\label{sec:bond-accrual}
Let \(d_s\) be settlement, \(d_M\) maturity, \(d_-\) the coupon date on or
before settlement and \(d_+\) the first coupon date after settlement.
The operator \(\operatorname{days}\) counts calendar days. The accrued
fraction \(\alpha\), remaining fraction \(w\), currency coupon \(C\) and
currency accrued interest \(AI\) satisfy
\begin{equation}\label{eq:accrual}
C=\frac{Fc}{m},\qquad
\alpha=\frac{\operatorname{days}(d_s-d_-)}
{\operatorname{days}(d_+-d_-)},\qquad w=1-\alpha,\qquad AI=C\alpha.
\end{equation}
% Explain maturity minus integer multiples of 12/m calendar months.
% At settlement on coupon date: alpha=0, w=1; exclude that day's coupon.
\section{Fractional-period YTM valuation}\label{sec:ytm-price}
For payment index \(i=1,\ldots,N\), \(N\) is the number of future payments
and \(d_i\) their dates. Define the fractional coupon-period count
\(q_i=w+i-1\), the coupon-based time \(T_i=q_i/m\) in years, and the currency
cash flow \(CF_i=C+F\mathbf{1}_{\{i=N\}}\).
The indicator \(\mathbf{1}_{\{i=N\}}\) equals one only at redemption.
The YTM-discounted present value is \(PV_i=CF_iB^{-q_i}\).
\begin{equation}\label{eq:ytm-price}
P_{\mathrm{dirty}}(\yytm)=\sum_{i=1}^{N}PV_i,\qquad
P_{\mathrm{clean}}=P_{\mathrm{dirty}}-AI.
\end{equation}
Both prices are currency amounts for \(F\); they are not automatically
normalised to 100. Hereafter \(P(\yytm)\) abbreviates the dirty YTM price.
% Explain fractional powers as an explicit discounting convention.
\section{Par, premium and discount relationships}\label{sec:par}
% Define premium/discount relative to face. Prove par only for coupon-date
% settlement with c=y_YTM, including the zero-yield case.
% Use monotonicity for premium/discount at that boundary. Do not assert
% dirty par or exactly clean par between coupons merely because c=y_YTM.
\plan{State the conditional par identity; move the geometric-series proof
to \cref{sec:app-par}. Place conventional context alongside
\textcite{tuckman2011}.}
\section{Duration and the analytical first derivative}\label{sec:duration}
% Define duration before its two forms; Macaulay is weighted time,
% modified is local proportional price sensitivity, not final maturity.
% Historical attribution: Macaulay (1938); textbook for contemporary conventions.
Macaulay duration \(D_{\mathrm{Mac}}\) is the present-value-weighted payment
time in years. Modified duration \(D_{\mathrm{Mod}}\) is the negative
dirty-price-normalised first derivative with respect to annual decimal YTM
\parencite{macaulay1938,tuckman2011}.
Primes on \(P\) denote derivatives with respect to \(\yytm\).
\begin{align}
P'(\yytm)&=-\sum_i CF_i\frac{q_i}{m}B^{-q_i-1},\label{eq:price-first}\\
D_{\mathrm{Mac}}&=\frac{\sum_i T_iPV_i}{P},&
D_{\mathrm{Mod}}&=-\frac{P'}P=\frac{D_{\mathrm{Mac}}}{B}.
\label{eq:duration}
\end{align}
% Explain fixed settlement and cash flows during differentiation, dirty-price
% normalisation, conventional year units and the role of m.
\section{Convexity and the analytical second derivative}\label{sec:convexity}
Convexity \(\mathcal C\) measures the curvature of the price--yield relation
through its normalised second derivative, conventionally in years squared.
\begin{equation}\label{eq:convexity}
P''(\yytm)=\sum_i CF_i\frac{q_i(q_i+1)}{m^2}B^{-q_i-2},
\qquad \mathcal C=\frac{P''}{P}.
\end{equation}
% Explain nonnegative cash flows imply positive convexity here.
% No factor 1/2 in convexity; it enters the Taylor expansion.
\section{DV01 and the Taylor approximation}\label{sec:ytm-dv01}
The decimal bump \(h\) equals \(10^{-4}\) for DV01, the currency sensitivity
for a one-basis-point move. The project uses down-minus-up central repricing:
\begin{equation}\label{eq:ytm-dv01}
\mathrm{DV01}_{\mathrm{YTM}}=
\frac{P(\yytm-h)-P(\yytm+h)}{2}
=hPD_{\mathrm{Mod}}+O(h^3).
\end{equation}
Here \(O(h^3)\) denotes a remainder bounded by a constant times \(|h|^3\)
as \(h\) approaches zero with other inputs fixed.
For a small annual decimal yield change \(\delta y\), Taylor expansion gives
\begin{equation}\label{eq:taylor-price}
\frac{P(\yytm+\delta y)-P(\yytm)}{P(\yytm)}
=-D_{\mathrm{Mod}}\delta y+\tfrac12\mathcal C(\delta y)^2
+O((\delta y)^3).
\end{equation}
% Define Taylor approximation as local polynomial expansion.
% DV01 is not a derivative divided by h. Both bumped yields must remain valid.
\section{Limitations of the selected bond conventions}
% No yield solver, irregular coupons, holiday adjustment, ex-coupon or credit model.
% Separate clean-price derivative equality from different price normalisations.
\plan{Link the limits of these identities to \cref{ch:limitations};
place detailed first/second differentiation in \cref{sec:app-derivatives}.}
